\documentclass[11pt,a4paper]{article}
\usepackage[utf8]{inputenc}
\usepackage[T1]{fontenc}
\usepackage{lmodern}
\usepackage[french]{babel}
\usepackage{amsmath,amssymb,mathtools}
\usepackage{icomma}
\usepackage{xcolor}
\usepackage{tikz}
\usepackage{pgfplots}
\usepackage[most]{tcolorbox}
\usepackage{enumitem}
\usepackage{fancyhdr}
\usepackage[hidelinks]{hyperref}
\usepackage{titlesec}
\usepackage[a4paper,margin=2.2cm,headheight=26pt,headsep=10pt]{geometry}
\usetikzlibrary{arrows.meta,positioning,calc,intersections,angles,quotes}
\usepgfplotslibrary{fillbetween}
\pgfplotsset{compat=1.18}
\definecolor{primary}{RGB}{30,75,145}
\definecolor{primarydark}{RGB}{20,50,100}
\definecolor{defcol}{RGB}{0,114,178}
\definecolor{thmcol}{RGB}{0,135,110}
\definecolor{excol}{RGB}{0,130,85}
\definecolor{remarkcol}{RGB}{120,70,180}
\definecolor{attncol}{RGB}{200,40,55}
\definecolor{methcol}{RGB}{85,85,100}
\definecolor{areacol}{RGB}{120,160,230}
\definecolor{areacol2}{RGB}{230,150,80}
\newcommand{\dd}{\mathrm{d}}
\newcommand{\R}{\mathbb{R}}
\newcommand{\ua}{\,\text{u.a.}}
\newcommand{\tg}{\tan}\newcommand{\cotg}{\cot}
\tcbset{boxbase/.style={enhanced,breakable,boxrule=0.9pt,arc=2.5pt,left=3mm,right=3mm,top=2.3mm,bottom=2.3mm,fonttitle=\bfseries,coltitle=white,toptitle=0.6mm,bottomtitle=0.6mm}}
\newtcolorbox{methode}[1][]{boxbase,colback=methcol!7,colframe=methcol,sharp corners,title={\textsc{Comment faire}\ifx\relax#1\relax\relax\else\textmd{ : \itshape #1}\fi}}
\newtcolorbox{exemple}[1][]{boxbase,colback=excol!5,colframe=excol,title={\textsc{Exemple}\ifx\relax#1\relax\relax\else\textmd{ : \itshape #1}\fi}}
\newtcolorbox{idee}[1][]{boxbase,colback=defcol!5,colframe=defcol,title={\textsc{Idée clé}\ifx\relax#1\relax\relax\else\textmd{ : \itshape #1}\fi}}
\newtcolorbox{remarque}[1][]{boxbase,colback=remarkcol!6,colframe=remarkcol,title={\textsc{Remarque}\ifx\relax#1\relax\relax\else\textmd{ : \itshape #1}\fi}}
\newtcolorbox{attention}[1][]{boxbase,colback=attncol!6,colframe=attncol,title={\textsc{Attention}\ifx\relax#1\relax\relax\else\textmd{ : \itshape #1}\fi}}
\newtcolorbox{exo}[1][]{boxbase,colback={primary!4},colframe=primary,title={\textsc{Exercice #1}}}
\newtcolorbox{corr}[1][]{boxbase,colback={thmcol!5},colframe=thmcol,title={\textsc{Corrigé #1}}}
\tikzset{axe/.style={->,>=Stealth,black!65,line width=0.7pt},courbe/.style={primary,line width=1.3pt},courbe2/.style={areacol2!90!black,line width=1.3pt},dvert/.style={gray,dashed,line width=0.7pt}}
\pagestyle{fancy}\fancyhf{}
\fancyhead[L]{\small\textcolor{primarydark}{\textbf{Fiche 7} — Intégrale d'une fonction continue}}
\fancyhead[R]{\small\textcolor{primarydark}{\textsc{Thème 4 — Intégrale définie}}}
\fancyfoot[C]{\small\thepage}
\renewcommand{\headrulewidth}{0.8pt}
\titleformat{\section}{\large\bfseries\color{primarydark}}{\thesection}{0.6em}{}

\begin{document}

\section*{Niveau Facile — Newton--Leibniz : de la primitive au nombre}

\begin{center}
\itshape\small Calculer les intégrales définies suivantes en appliquant la formule de
Newton--Leibniz $\int_a^b f=[F(x)]_a^b=F(b)-F(a)$. Les primitives sont toutes immédiates.
\end{center}

\medskip

\begin{exo}[1 — polynôme, forme de base]
Calculer
\[
I=\int_0^1\bigl(3x^2+2x+1\bigr)\,\dd x.
\]
Détailler la primitive puis l'évaluation par la notation crochet.
\end{exo}

\begin{exo}[2 — bornes négatives]
Calculer
\[
J=\int_{-1}^{2}\bigl(x^2-x\bigr)\,\dd x.
\]
Soigner le calcul de $F(-1)$ (attention aux signes).
\end{exo}

\begin{exo}[3 — bien poser le crochet]
Calculer
\[
K=\int_{1}^{3}\bigl(2x-1\bigr)\,\dd x
\]
en écrivant explicitement la primitive entre crochets $\bigl[\,F(x)\,\bigr]_1^3$ avant de
substituer.
\end{exo}

\begin{exo}[4 — intégrale nulle et inversion des bornes]
\begin{enumerate}[label=\alph*),leftmargin=1.8em]
  \item Que vaut $\displaystyle\int_{5}^{5}\bigl(x^3+7\bigr)\,\dd x$\,? Justifier sans calcul.
  \item Calculer $\displaystyle L=\int_{1}^{3} x\,\dd x$, puis en déduire \emph{sans nouveau calcul}
        la valeur de $\displaystyle\int_{3}^{1} x\,\dd x$.
\end{enumerate}
\end{exo}

\begin{exo}[5 — une puissance fractionnaire]
Calculer
\[
M=\int_{1}^{4}\sqrt{x}\,\dd x.
\]
\emph{Indication} : $\sqrt{x}=x^{1/2}$, donc on utilise la forme puissance avec $n=\tfrac12$.
\end{exo}

\end{document}
